%% kapa: exam=1 round=35 year=2024 period=1 form=A subject=civil_law subject_ko=민법 seq=30 number=30
%% source: Q-Net 감정평가사 기출문제 (cst00309) artlSeq=5213229 file=raw/1008/5213229/1교시 문제지 원본.pdf
%% conversion: kordoc 4.24.51 -> tools/gen1.py (auto); review=proofread-sonnet; flags=none
\begin{question}{30}
甲은 그 소유 X토지에 대한 배타적 사용ㆍ수익권을 포기하고 타인(사인, 국가 등 일반 공중)의 통행을 위한 용도로 제공하였다. 이에 관한 설명으로 옳지 않은 것은? (다툼이 있으면 판례에 따름)
\begin{choices}
  \item 甲은 그 타인에 대하여 X의 인도청구를 할 수 없다.
  \item 甲이 X에 대한 소유권을 보유한 채 사용ㆍ수익권을 대세적ㆍ영구적으로 포기하는 것은 허용되지 않는다.
  \item 甲은 일반 공중의 통행을 방해하지 않는 범위에서 X를 처분할 수 있다.
  \item 甲의 상속인의 X에 대한 배타적 사용ㆍ수익권도 제한된다.
  \item 만약 甲이 X를 일반 공중의 통행목적이 아니라 지상건물의 소유자만을 위하여 배타적 사용ㆍ수익권을 포기한 경우, 특별한 사정이 없는 한 X의 매수인의 배타적 사용ㆍ수익권 행사는 제한된다.
\end{choices}
\end{question}
\answer{5}
